\section{Random SPIN}
\label{sec:random-spin}

Random SPIN replaces the accumulator with a random recursive convolution.
Its outer uses the independent random-block injections from
Section~\ref{sec:accumulator-spin}.

\subsection{Relation to Silver and Expand-Convolute}

Silver codes predate Expand-Convolute and use a banded lower-triangular solve
for linear-time encoding~\cite{CouteauRindalRaghuraman2021Silver}. The solve
is an earlier structured recursive encoder of convolutional form. Silver does
not sample independent convolution taps at every time step.

Expand-Convolute makes the random convolution explicit. It combines a sparse
expander with a non-wrapping recursion whose time-dependent feedback
coefficients are independent uniform bits
\cite{RaghuramanRindalTanguy2023EC}. The Random SPIN theorem retains that
inner distribution and replaces the expander with independent random-block
injections followed by one uniform interleaver.

\subsection{Random recursive inner}

Fix a memory $m\ge1$. For every $t\in[N]$, setup samples
\[
  \alpha_t\gets\F_2^m
\]
independently. Set $y_j:=0$ for $j\le0$. On input $u\in\F_2^N$, define
\begin{equation}
  y_t
  :=
  u_t+
  \left\langle
    \alpha_t,(y_{t-1},\ldots,y_{t-m})
  \right\rangle,
  \qquad t\in[N].
  \label{eq:random-recursive-inner}
\end{equation}
Write $I_{N,m}^{\mathrm{rnd}}(u):=y$. Once setup fixes
$(\alpha_t)_{t=1}^N$, this map is linear and invertible because $y_t$ contains
$u_t$ with coefficient one.

\subsection{Exact convolution transfer}

The random inner admits an exact finite-state transform. Let $X$ mark input
weight and let $S$ mark output weight. Index the matrix rows and columns by
$z\in\{0,\ldots,m\}$, where $z$ is the number of trailing output zeros,
clipped at $m$. Thus $z=m$ means that the convolution state is zero.

Define the nonnegative matrix $T_m(X,S)$ by
\begin{equation}
  T_m(X,S)[z,z']
  :=
  \begin{cases}
    (1+X)S/2, & z<m,\ z'=0,\\
    (1+X)/2, & z<m,\ z'=z+1,\\
    XS, & z=m,\ z'=0,\\
    1, & z=m,\ z'=m,\\
    0, & \text{otherwise}.
  \end{cases}
  \label{eq:random-transfer-matrix}
\end{equation}
Let $e_m$ be the unit vector at index $m$, and let $\boldsymbol 1$ be the
all-ones vector.

\begin{lemma}[Exact random-convolution transform]
\label{lem:exact-random-convolution-transform}
For the setup in~\eqref{eq:random-recursive-inner},
\begin{equation}
  \E\!\left[
    \sum_{u\in\F_2^N}
    X^{\wt(u)}S^{\wt(I_{N,m}^{\mathrm{rnd}}(u))}
  \right]
  =
  e_m^{\mathsf T}T_m(X,S)^N\boldsymbol 1.
  \label{eq:exact-random-convolution-transform}
\end{equation}
\end{lemma}

\begin{proof}
Suppose $z<m$. The convolution state is nonzero, so the fresh uniform
$\alpha_t$ makes the output bit uniform for either input bit. Summing the two
input choices with weights $1$ and $X$ gives the first two nonzero transitions
in~\eqref{eq:random-transfer-matrix}. If $z=m$, input zero stays in the zero
state, while input one produces output one and moves to state zero. These are
the remaining two transitions. Multiplying the one-step transforms and
summing over the terminal state proves the identity.
\end{proof}

Let $U_h$ be uniform on the weight-$h$ slice, independently of the inner
setup. For $0<S<1$, $a>0$, and $D\in\{0,\ldots,N\}$,
Lemma~\ref{lem:exact-random-convolution-transform} gives
\begin{equation}
  \Prob[\wt(I_{N,m}^{\mathrm{rnd}}(U_h))\le D]
  \le
  S^{-D}a^{-h}
  \frac{e_m^{\mathsf T}T_m(a,S)^N\boldsymbol 1}{\binom Nh}.
  \label{eq:random-convolution-tail}
\end{equation}
The first step is the lower-tail moment bound. The second bounds the
$X^h$ coefficient by evaluating the nonnegative polynomial at $a$.

For $\delta\in(0,1/2)$, define
\begin{equation}
  p_{N,h}(\delta)
  :=
  \Prob[
    \wt(I_{N,m}^{\mathrm{rnd}}(U_h))
    \le\lfloor\delta N\rfloor
  ].
  \label{eq:random-inner-slice-tail}
\end{equation}

\subsection{Asymptotic theorem}

For an even block length $B$ dividing $N$, set $M:=N/B$. Independently
sample uniform linear injections
\[
  C_i:\F_2^{B/2}\longrightarrow\F_2^B,
  \qquad i\in[M],
\]
one permutation $\Pi\gets S_N$, and the convolution setup above. Define
\begin{equation}
  E_{N,m}^{\mathrm{rnd}}(x)
  :=
  I_{N,m}^{\mathrm{rnd}}\!\left(
    \Pi(C_1(x^{(1)})\Vert\cdots\Vert C_M(x^{(M)}))
  \right).
  \label{eq:independent-random-spin}
\end{equation}
All setup objects are independent and remain fixed for all messages.

For each sufficiently large $M$, let $B_M^*$ be the least positive even
integer satisfying
\[
  B_M^*\ge17\log_2(MB_M^*).
\]
Set
\[
  N_M^*:=MB_M^*,
  \qquad
  m_M^*:=\left\lceil\frac{51}{50}\log_2N_M^*\right\rceil,
  \qquad
  \delta_*:=\frac{5501}{50000}.
\]

\begin{theorem}[Random SPIN]
\label{thm:independent-random-spin}
For the ensemble in~\eqref{eq:independent-random-spin},
\begin{equation}
  \Prob[
    d_{\min}(E_{N_M^*,m_M^*}^{\mathrm{rnd}})
    \le\delta_*N_M^*
  ]
  =o(1)
  \qquad(M\to\infty).
  \label{eq:independent-random-spin-distance}
\end{equation}
The family has rate $1/2$, relative distance greater than $0.11002$,\linebreak and
$\Theta(N_M^*\log N_M^*)$ direct encoding cost.
\end{theorem}

\begin{proof}
Let $A_{N,h}^{\mathrm{out}}$ be the number of outer words of weight $h$, and
let $p_{N,h}(\delta_*)$ be the probability in
\eqref{eq:random-convolution-tail} with
$D=\lfloor\delta_*N\rfloor$. Split the first-moment sum at
$h=N/5000$ and $h=0.9N$.

Lemma~\ref{lem:random-spin-sparse-closure} shows that the first range is
$o(1)$. Lemma~\ref{lem:random-spin-linear-closure} is uniform for
$h/N\in[1/5000,0.9]$, so the middle range is $2^{-\Omega(N)}$. For the final
range, use $p_{N,h}\le1$ and the rate-half random-block outer exponent:
\[
  \sum_{h\ge0.9N}\E[A_{N,h}^{\mathrm{out}}]
  \le
  2^{N(H_2(0.9)-1/2+o(1))}
  =2^{-\Omega(N)}.
\]
The complete first moment tends to zero. Markov's inequality proves
\eqref{eq:independent-random-spin-distance}.
\end{proof}

The admissible lengths have relative gaps $o(1)$. For a requested length, zero
extend the largest preceding native member. This preserves injectivity and
absolute minimum distance while changing rate and relative distance by only
$o(1)$.
